\documentclass[11pt,a4paper]{article}

\usepackage[utf8]{inputenc}
\usepackage[T1]{fontenc}
\usepackage{geometry}
\usepackage{hyperref}
\usepackage{enumitem}
\usepackage{xcolor}
\usepackage{fancyhdr}
\usepackage{amsmath}
\usepackage{amssymb}
\usepackage{booktabs}
\usepackage{listings}

\geometry{margin=1in}
\hypersetup{
  colorlinks=true,
  linkcolor=blue,
  urlcolor=blue
}

\setlist{itemsep=0.35em, topsep=0.35em}

\pagestyle{fancy}
\fancyhf{}
\lhead{Elements of AIML (B.Tech)}
\rhead{Unit 01: Diagnostic Assessment}
\cfoot{\thepage}
\setlength{\headheight}{14pt}

\lstset{
  basicstyle=\ttfamily\small,
  keywordstyle=\color{blue},
  commentstyle=\color{gray},
  stringstyle=\color{teal},
  showstringspaces=false,
  columns=fullflexible,
  frame=single
}

\newcommand{\SubmissionDeadline}{March 8, 2026}

% Marks per question (total = 100)
\newcommand{\EMarks}{\textbf{[2 marks]} }
\newcommand{\MMarks}{\textbf{[3 marks]} }
\newcommand{\HMarks}{\textbf{[5 marks]} }

\begin{document}

\begin{center}
  {\LARGE \textbf{Diagnostic Assessment -- Unit 01}}\\[0.25em]
  {\large Elements of AIML}\\[0.25em]
  \normalsize (Introduction to AI; designed to identify learning gaps)
\end{center}

\noindent
\textbf{Instructor:} Tofik Ali \hfill \textbf{Time Limit:} None\\
\textbf{Submission Deadline:} \SubmissionDeadline\\
\textbf{Student Name:} \_\_\_\_\_\_\_\_\_\_\_\_\_\_\_\_\_ \hfill \textbf{Roll No.:} \_\_\_\_\_\_\_\_\\

\vspace{0.6em}
\hrule
\vspace{0.8em}

\textbf{Instructions}
\begin{itemize}[leftmargin=*]
  \item \textbf{This is a diagnostic assessment, not a quiz/test.} The goal is to identify learning gaps so you know what to revise next.
  \item \textbf{Total marks: 100.} Easy: 10 $\times$ 2 = 20, Medium: 10 $\times$ 3 = 30, Hard: 10 $\times$ 5 = 50.
  \item \textbf{No time limit.} Submit on or before \textbf{\SubmissionDeadline}.
  \item \textbf{Important (70\% rule): To have your assignment marks counted, you must score at least 70\% (70/100 or more) in this assessment.}
    If you score below 70\%, you may reattempt and resubmit. Your latest submission will be considered.
  \item \textbf{Marking policy (honest attempt):} Marks are deducted only for (i) not attempting a question and (ii) high similarity with other submissions (copying).
    Correctness is \textbf{not} the main focus; your reasoning and effort matter more.
  \item \textbf{Show your real effort:} if you tried multiple approaches for any question, write them all (Attempt 1, Attempt 2, ...). Partial/incorrect attempts are acceptable.
  \item \textbf{Many questions are open-ended by design:} multiple solutions/variants are acceptable. Use your own example data and explain your choices.
  \item Unless specified, you may choose your own symbols/assumptions (e.g., how you name predicates), but you must clearly mention your choices.
  \item Submission format: submit a single ZIP containing \texttt{notebook.ipynb} and \texttt{README.md} (plus optional \texttt{assets/}).
  \item Round final numeric answers to \textbf{2 decimal places} (when applicable).
\end{itemize}

\vspace{0.6em}

\section*{Easy (10 Questions)}
\begin{enumerate}[label=\textbf{E\arabic*.}, leftmargin=*]
  \item \EMarks \textbf{[Subjective]} Define AI in one sentence and give one example application.

  \item \EMarks \textbf{[MCQ -- Select all that apply]} Which of the following are \textbf{AI tasks} (as commonly defined in Unit~01)?
    \begin{enumerate}[label=(\Alph*), leftmargin=*]
      \item Sorting numbers in ascending order
      \item Face recognition in a photo
      \item Route planning with changing traffic
      \item A calculator computing $23 \times 47$
      \item Spam detection in emails
    \end{enumerate}

  \item \EMarks \textbf{[Subjective]} Write 2--3 lines each: AI vs ML vs Data Science.

  \item \EMarks \textbf{[MCQ -- Select all that apply]} Which items belong to PEAS?
    \begin{enumerate}[label=(\Alph*), leftmargin=*]
      \item Performance measure
      \item Environment
      \item Algorithm
      \item Actuators
      \item Sensors
    \end{enumerate}

  \item \EMarks \textbf{[Subjective]} Create a PEAS description for a \textbf{smart vacuum cleaner} (2--4 bullets per PEAS item).

  \item \EMarks \textbf{[MCQ -- Select one]} Chess is best described as:
    \begin{enumerate}[label=(\Alph*), leftmargin=*]
      \item Fully observable and deterministic
      \item Partially observable and stochastic
      \item Fully observable and stochastic
      \item Partially observable and deterministic
    \end{enumerate}

  \item \EMarks \textbf{[Subjective]} What is a \textbf{rational agent}? (3--5 lines)

  \item \EMarks \textbf{[MCQ -- Select all that apply]} Which environments are typically \textbf{dynamic}?
    \begin{enumerate}[label=(\Alph*), leftmargin=*]
      \item Playing chess (during your move)
      \item Driving a car in traffic
      \item Solving a static Sudoku puzzle
      \item Stock trading during market hours
    \end{enumerate}

  \item \EMarks \textbf{[Subjective]} List the four main components of a \textbf{learning agent}.
    Write one line about the role of the \textbf{critic}.

  \item \EMarks \textbf{[Subjective]} List two limitations of AI in real deployments and one practical mitigation for each.
\end{enumerate}

\section*{Medium (10 Questions)}
\begin{enumerate}[label=\textbf{M\arabic*.}, leftmargin=*]
  \item \MMarks \textbf{[Subjective]} Create a PEAS description for a \textbf{navigation app} that suggests routes on campus.

  \item \MMarks \textbf{[Subjective]} Classify the environment of an \textbf{online shopping recommender} using at least five properties
    (observable, deterministic, episodic, static, discrete, etc.). Justify briefly.

  \item \MMarks \textbf{[Subjective]} Choose an agent type (simple reflex / model-based / goal-based / utility-based / learning)
    for a \textbf{smart thermostat}. Justify in 6--8 lines.

  \item \MMarks \textbf{[Subjective]} For a \textbf{fraud detection} system, write 3 success criteria beyond accuracy
    (e.g., cost, safety, fairness, false alarms). Explain in 6--8 lines.

  \item \MMarks \textbf{[Subjective]} Explain performance measure vs reward with a concrete example (6--8 lines).

  \item \MMarks \textbf{[Subjective]} Give one realistic example of bias in an AI system.
    Mention the data source and one mitigation (8--10 lines).

  \item \MMarks \textbf{[Subjective]} Give one example of distribution shift that can break an AI system.
    Propose 2 monitoring signals you would track after deployment.

  \item \MMarks \textbf{[Subjective]} Make a table with 3 technique categories (search/planning, knowledge/logic, learning)
    and 2 examples for each.

  \item \MMarks \textbf{[Subjective]} Consider an AI-based proctoring tool.
    List 3 risks (privacy/fairness/safety) and one control for each.

  \item \MMarks \textbf{[Subjective]} ``Model accuracy can be high but system success can be low.'' Explain with one example (8--10 lines).
\end{enumerate}

\section*{Hard (10 Questions)}
\begin{enumerate}[label=\textbf{H\arabic*.}, leftmargin=*]
  \item \HMarks \textbf{[Subjective]} Design an AI agent for an \textbf{autonomous delivery robot} inside a campus.
    Provide PEAS, environment properties, and success criteria.

  \item \HMarks \textbf{[Subjective]} A face-recognition attendance system performs poorly for some groups.
    Write a short analysis: what could be the causes, how to measure the issue, and how to improve it.

  \item \HMarks \textbf{[Subjective]} Write an 8--12 bullet deployment checklist for a student-facing AI feature
    (privacy, monitoring, fallback, human override, logging, etc.).

  \item \HMarks \textbf{[Subjective]} Compare rule-based vs learning-based approaches for two tasks:
    (i) spam filtering, (ii) timetable scheduling. Explain trade-offs (10--12 lines).

  \item \HMarks \textbf{[Subjective]} Write a 1-page Responsible AI note for a university AI tool
    including fairness, privacy, transparency, accountability, and safety.

  \item \HMarks \textbf{[Subjective]} Explain why self-driving is difficult using environment properties
    (at least 6 properties with justification).

  \item \HMarks \textbf{[Subjective]} Describe the learning agent components for a recommendation system.
    Draw a block diagram and explain the feedback loop.

  \item \HMarks \textbf{[Subjective]} Choose one AI application where explainability is necessary.
    Explain what kind of explanation is needed and who needs it.

  \item \HMarks \textbf{[Subjective]} Write a short impact analysis (1 page) of AI in one domain:
    education / healthcare / banking / security. Include both benefits and harms.

  \item \HMarks \textbf{[Subjective]} Write a short reflection (12--16 lines):
    ``Where is AI overhyped today, and where is it genuinely useful?'' Support with examples.
\end{enumerate}

\vfill
\noindent\textit{End of Assessment}

\end{document}
